% =============================================================================
% demo-beautiful.tex — Demo document for the swarmbeauty theme
%
% Compile: lualatex --shell-escape demo-beautiful.tex
%          (run twice for TOC)
% =============================================================================

\documentclass[11pt, a4paper]{scrartcl}

\usepackage{swarmbeauty}

% ── Metadata ─────────────────────────────────────────────────────────────────

\title{SwarmBeauty Theme Demo}
\subtitle{A showcase of the all-in-one LaTeX Beautiful Theme}
\author{LaTeX Helper Swarm Project}
\date{May 14, 2026}

% ── Bibliography (optional — uncomment if you have a .bib file) ─────────────
% \usepackage[backend=biber, style=numeric]{biblatex}
% \addbibresource{references.bib}

\begin{document}

% ══════════════════════════════════════════════════════════════════════════════
% Title Page
% ══════════════════════════════════════════════════════════════════════════════

\swarmtitlepage

% ══════════════════════════════════════════════════════════════════════════════
% Table of Contents
% ══════════════════════════════════════════════════════════════════════════════

\tableofcontents
\newpage

% ══════════════════════════════════════════════════════════════════════════════
% Section 1: Introduction
% ══════════════════════════════════════════════════════════════════════════════

\section{Introduction}

This document demonstrates the capabilities of the \code{swarmbeauty} theme, a comprehensive
LaTeX document style designed for clarity, elegance, and usability. The theme is built on
top of KOMA-Script and provides a rich set of features out of the box, including beautiful
section headings, styled tables, syntax-highlighted code blocks, and custom block
environments for notes, tips, and warnings.

The design philosophy behind this theme prioritizes \emphtext{readability} and
\emphtext{professional aesthetics} while maintaining compatibility with the standard
LaTeX ecosystem. Every element---from the title page to the footnote---has been
carefully crafted to produce documents that look polished and modern.

\subsection{Motivation}

Writing beautiful LaTeX documents typically requires dozens of package inclusions,
meticulous configuration of colors, fonts, and spacing, and a deep understanding of
the underlying typesetting system. The \code{swarmbeauty} theme aims to eliminate this
friction by providing sensible defaults that work well for academic papers, technical
reports, and documentation.

\begin{noteblock}
This theme requires \textbf{LuaLaTeX} or \textbf{XeLaTeX} with the
\code{-{}-shell-escape} flag for the minted code highlighting.
\end{noteblock}

% ══════════════════════════════════════════════════════════════════════════════
% Section 2: Typography & Text
% ══════════════════════════════════════════════════════════════════════════════

\section{Typography and Text}

The theme uses \textbf{Latin Modern} fonts throughout, with \textit{italic},
\textbf{bold}, and \textbf{\textit{bold italic}} variants. Mathematics is rendered
using \texttt{unicode-math} for consistent glyph quality. Microtypographic
improvements via \code{microtype} provide better character spacing and hyphenation.

\subsection{Inline Elements}

You can highlight \hltext{important terms} using the \code{\\hltext} command.
For inline code, use the \code{\\code} command like this: \code{some_function()}.
Emphasis boxes are available too: \emphtext{key concept}.

\subsection{Lists}

\begin{itemize}
  \item First item in an unordered list
  \item Second item with \emphtext{emphasis}
  \item Third item with nested content:
    \begin{itemize}
      \item Nested item A
      \item Nested item B
    \end{itemize}
\end{itemize}

\begin{enumerate}
  \item Step one: read the documentation
  \item Step two: include the theme
  \item Step three: write beautiful documents
\end{enumerate}

% ══════════════════════════════════════════════════════════════════════════════
% Section 3: Block Environments
% ══════════════════════════════════════════════════════════════════════════════

\section{Block Environments}

The theme provides five distinct block environments, each with a unique color
accent and left border for quick visual identification.

\begin{noteblock}
This is a \textbf{note} block. Use it for general observations, additional
context, or supplementary information that is relevant but not critical to
the main argument.
\end{noteblock}

\begin{tipblock}
This is a \textbf{tip} block. Use it to highlight best practices, shortcuts,
or recommendations that can improve the reader's workflow or understanding.
\end{tipblock}

\begin{warningblock}
This is a \textbf{warning} block. Use it to alert the reader about potential
pitfalls, common mistakes, or deprecated features that they should be aware of.
\end{warningblock}

\begin{dangerblock}
This is a \textbf{danger} block. Use it for critical warnings about data loss,
security vulnerabilities, or operations that cannot be undone.
\end{dangerblock}

\begin{exampleblock}
This is an \textbf{example} block. Use it to present worked examples, sample
output, or illustrations of the concepts being discussed.
\end{exampleblock}

% ══════════════════════════════════════════════════════════════════════════════
% Section 4: Tables
% ══════════════════════════════════════════════════════════════════════════════

\section{Tables}

The theme includes custom table rules and integrates seamlessly with both
\code{booktabs} and \code{tabularray}. Below are examples of each style.

\subsection{Booktabs Table}

\begin{table}[H]
  \centering
  \caption{Comparison of LaTeX compilation engines}
  \label{tab:engines}
  \begin{tabular}{l c c c}
    \swarmtoprule
    \textbf{Engine}   & \textbf{Speed} & \textbf{Font Support} & \textbf{Unicode} \\
    \swarmmidrule
    pdfLaTeX           & Fast           & Limited (Type 1)      & Partial          \\
    XeLaTeX            & Medium         & Full (OpenType)        & Full             \\
    LuaLaTeX           & Slow           & Full (OpenType)        & Full             \\
    \swarmbottomrule
  \end{tabular}
\end{table}

\subsection{Tabularray Table}

\begin{table}[H]
  \centering
  \caption{Key packages used by the theme}
  \label{tab:packages}
  \begin{tblr}{
    colspec     = {l X[c] X[c]},
    row{1}      = {font=\sffamily\bfseries, bg=sbLight},
    row{odd}    = {bg=sbCodeBg},
    hline{1}    = {1pt, sbPrimary},
    hline{2}    = {0.4pt, sbMedium},
    hline{Z}    = {1pt, sbPrimary},
    vlines,
  }
    Package       & Purpose                              & Required \\
    fontspec      & OpenType font loading                & Yes      \\
    microtype     & Microtypographic improvements        & Yes      \\
    minted        & Pygments-based code highlighting     & Yes      \\
    tcolorbox     & Styled boxes and code blocks          & Yes      \\
    booktabs      & Professional table rules             & Yes      \\
    tabularray    & Modern table formatting              & Optional \\
    tikz          & Graphics and decorations             & Yes      \\
  \end{tblr}
\end{table}

% ══════════════════════════════════════════════════════════════════════════════
% Section 5: Code Listings
% ══════════════════════════════════════════════════════════════════════════════

\section{Code Listings}

The \code{codeblock} environment provides syntax-highlighted code with line
numbers, powered by the Pygments library through the \code{minted} package.

\begin{codeblock}[Python]{Fibonacci sequence generator}
def fibonacci(n: int) -> list[int]:
    """Generate the first n Fibonacci numbers."""
    seq = [0, 1]
    for i in range(2, n):
        seq.append(seq[-1] + seq[-2])
    return seq[:n]

if __name__ == "__main__":
    result = fibonacci(10)
    print(f"First 10: {result}")
\end{codeblock}

\begin{codeblock}[LaTeX]{A simple equation}
\begin{equation}
  E = mc^2
  \label{eq:einstein}
\end{equation}
\end{codeblock}

% ══════════════════════════════════════════════════════════════════════════════
% Section 6: Mathematics
% ══════════════════════════════════════════════════════════════════════════════

\section{Mathematics}

The theme provides styled theorem environments with colored left borders.

\begin{theorem}{Euler's Identity}{euler}
  For any real number $x$, the following identity holds:
  \[
    e^{i\pi} + 1 = 0
  \]
\end{theorem}

\begin{definition}{Metric Space}{metric}
  A \emph{metric space} is a pair $(X, d)$ where $X$ is a set and $d: X \times X \to \mathbb{R}_{\geq 0}$ satisfies:
  \begin{enumerate}
    \item $d(x, y) = 0 \iff x = y$ (identity of indiscernibles)
    \item $d(x, y) = d(y, x)$ (symmetry)
    \item $d(x, z) \leq d(x, y) + d(y, z)$ (triangle inequality)
  \end{enumerate}
\end{definition}

\begin{lemma}{Nested Intervals}{nested}
  Let $\{[a_n, b_n]\}_{n=1}^{\infty}$ be a sequence of closed, bounded intervals such that
  $[a_{n+1}, b_{n+1}] \subseteq [a_n, b_n]$ for all $n$. Then
  $\bigcap_{n=1}^{\infty} [a_n, b_n] \neq \emptyset$.
\end{lemma}

\subsection{Inline and Display Math}

The quadratic formula: $x = \frac{-b \pm \sqrt{b^2 - 4ac}}{2a}$

A more complex display equation:
\[
  \int_{-\infty}^{\infty} e^{-x^2} \, dx = \sqrt{\pi}
\]

% ══════════════════════════════════════════════════════════════════════════════
% Section 7: Figures
% ══════════════════════════════════════════════════════════════════════════════

\section{Figures}

Figures use standard \code{figure} environments with the theme's caption styling.

\begin{figure}[H]
  \centering
  \begin{tikzpicture}
    \draw[thick, sbSecondary, fill=sbLight, rounded corners=4pt]
      (0,0) rectangle (6,3);
    \node[sfont, sbDark] at (3, 1.5) {Your figure goes here};
  \end{tikzpicture}
  \caption{A placeholder figure rendered with TikZ}
  \label{fig:placeholder}
\end{figure}

\begin{tipblock}
Replace the placeholder above with \code{\\includegraphics[width=\\linewidth]\{image.pdf\}}.
\end{tipblock}

% ══════════════════════════════════════════════════════════════════════════════
% Section 8: Summary
% ══════════════════════════════════════════════════════════════════════════════

\section{Summary}

This demo has showcased the following features of the \code{swarmbeauty} theme:

\begin{enumerate}
  \item \textbf{Title page} with colored header bar and metadata box
  \item \textbf{Section headings} with decorative rules
  \item \textbf{Block environments} (note, tip, warning, danger, example)
  \item \textbf{Styled tables} with booktabs and tabularray
  \item \textbf{Syntax-highlighted code} via minted + tcolorbox
  \item \textbf{Theorem environments} with colored borders
  \item \textbf{Typography} with Latin Modern fonts and microtype
  \item \textbf{Headers and footers} with page numbers and document info
\end{enumerate}

\colorrule

\begin{center}
  \small\textcolor{sbMedium}{%
    Generated with \texttt{swarmbeauty} theme --- LaTeX Helper Swarm Project%
  }
\end{center}

\end{document}
